% Source: https://artofproblemsolving.com/community/c3h1387713_2016_mathcounts_nationals_21
% Stable identity: aops-topic-1387713-q21
Treat the six socks as distinct and work day by day.

\textbf{Day 1.} There are $\dbinom62=15$ equally likely pairs, of which $3$ are matching pairs, so the first pair is mixed with probability $\dfrac{12}{15}=\dfrac{4}{5}$.

\textbf{Day 2.} Say the first pair was red and white. The remaining socks are one red, one white and two blue, giving $\dbinom42=6$ equally likely pairs: red-white ($1$ pair), red-blue ($2$), white-blue ($2$), blue-blue ($1$). The blue-blue pair fails on day 2, and the red-white pair leaves blue-blue for day 3, which also fails. The other $4$ pairs leave a mixed pair for day 3. So the probability of success from here is $\dfrac{4}{6}=\dfrac{2}{3}$.

Hence the probability is
\[
\dfrac{4}{5}\cdot\dfrac{2}{3}=\boxed{\dfrac{8}{15}} .
\]

\textbf{The error.} The student's $21$ ``cases'' are not equally likely: they lump together different numbers of actual sock choices (for instance there are $3$ matching first pairs but $12$ mixed first pairs, and the second-day choices also have different numbers of socks behind them). A probability can be found as a ratio of counts only when the counted outcomes are equally likely. Counting distinct socks, as above, fixes this.
